
\section*{Exorde : Le constat de l'effondrement (quantitatif et qualitatif)}

Le grec ancien est en train de disparaître. Non pas lentement, insidieusement, mais de manière spectaculaire et mesurable. En France, le nombre d'élèves étudiant le grec ancien a chuté de plus de 90\% en quarante ans. Là où existaient autrefois des milliers de lycéens maîtrisant tant bien que mal quelques passages de Platon ou d'Homère, ne restent aujourd'hui que quelques centaines d'une élite déjà fragmentée. Les universités ferment leurs départements de Lettres Classiques. Les agrégations en grec ne proposent plus que des postes de recherche, la formation pédagogique s'évapore.

Mais le constat n'est pas seulement statistique : il est qualitatif. Ceux qui persistent à apprendre le grec ancien le font de moins en moins bien. Les trois ou quatre années d'étude réglementaire au lycée produisent rarement un locuteur capable de lire avec aisance, sans dictionnaire, une page de texte ancien. L'objectif pédagogique historique---former des esprits cultivés, capables de converser avec les monuments de la pensée classique---s'est évaporé. À la place : une culture de la traduction laborieuse, de l'analyse grammaticale stérile, du déchiffrement sans compréhension vivante.

Le paradoxe est brutal : jamais le grec ancien n'a été aussi prestigieux culturellement. Les racines grecques dominent les vocabulaires scientifiques mondiaux. Le patrimoine antique a gagné en aura dans l'imaginaire collectif (mythologie dans les séries, jeux vidéo, littérature jeunesse). Et pourtant, l'accès direct à cette langue, l'instrument même de sa pensée, se dérobe.
